\documentclass[aspectratio=169,11pt]{beamer}

\usetheme{AnnArbor}
\usecolortheme{dolphin}
\usefonttheme{professionalfonts}
\setbeamertemplate{navigation symbols}{}
\setbeamertemplate{caption}[numbered]
\setbeamertemplate{itemize items}[circle]
\setbeamertemplate{enumerate items}[default]

\usepackage{fontspec}
\setsansfont{DejaVu Sans}
\setmonofont{DejaVu Sans Mono}[Scale=0.85]
\usepackage{amsmath,amssymb}
\usepackage{booktabs}
\usepackage{array}
\usepackage{tabularx}
\usepackage{adjustbox}
\usepackage{listings}
\usepackage{xcolor}
\usepackage{tikz}
\usetikzlibrary{arrows.meta,positioning,shapes.geometric,fit,calc}
\usepackage{hyperref}

% Mau sac thuong hieu DBMS
\definecolor{dbblue}{RGB}{20,74,135}
\definecolor{dblight}{RGB}{235,243,255}
\definecolor{dbgreen}{RGB}{28,120,90}
\definecolor{dborange}{RGB}{210,120,30}
\definecolor{dbred}{RGB}{185,28,28}
\definecolor{codebg}{RGB}{246,248,250}
\definecolor{codekw}{RGB}{0,80,180}
\definecolor{codestr}{RGB}{180,60,20}
\definecolor{codecomm}{RGB}{100,110,120}

\hypersetup{colorlinks=true,linkcolor=dbblue,urlcolor=dbblue}

\setbeamercolor{structure}{fg=dbblue}
\setbeamercolor{title}{fg=white,bg=dbblue}
\setbeamercolor{frametitle}{fg=white,bg=dbblue}
\setbeamercolor{block title}{fg=white,bg=dbblue}
\setbeamercolor{block body}{bg=dblight,fg=black}
\setbeamercolor{block title example}{fg=white,bg=dbgreen}
\setbeamercolor{block body example}{bg=dblight!60!white,fg=black}
\setbeamercolor{block title alerted}{fg=white,bg=dborange}
\setbeamercolor{block body alerted}{bg=dblight!40!white,fg=black}

\lstset{
  basicstyle=\ttfamily\scriptsize,
  keywordstyle=\color{codekw}\bfseries,
  stringstyle=\color{codestr},
  commentstyle=\color{codecomm}\itshape,
  backgroundcolor=\color{codebg},
  frame=single,
  rulecolor=\color{dbblue!30},
  breaklines=true,
  showstringspaces=false,
  keepspaces=true,
  aboveskip=2pt,
  belowskip=2pt,
  language=SQL,
  morekeywords={START,TRANSACTION,COMMIT,ROLLBACK,SAVEPOINT,RELEASE,AUTOCOMMIT,ISOLATION,LEVEL,READ,UNCOMMITTED,COMMITTED,REPEATABLE,SERIALIZABLE,LOCK,UNLOCK,TABLES,SHARED,EXCLUSIVE,ENGINE,INNODB,MYISAM},
  tabsize=2
}

\newcommand{\code}[1]{\texttt{#1}}

\title[Thao tác Dữ liệu \& Giao dịch MySQL]{Thao tác Dữ liệu và Giao dịch trong MySQL}
\subtitle{Transaction, Mô hình ACID, Kiểm soát Autocommit, Savepoints \& Mức độ Cô lập}
\author[Vũ Đức Minh]{Vũ Đức Minh}
\institute[NEU]{Khoa Công nghệ Thông tin --- Trường Đại học Kinh tế Quốc dân (NEU)}
\date{Năm học 2026--2027}

\begin{document}

% Slide 1: Trang tieu de
\begin{frame}
  \titlepage
\end{frame}

% Slide 2: Muc tieu bai hoc
\begin{frame}{Mục tiêu Học tập (Learning Objectives)}
  \begin{itemize}
    \item \textbf{Hiểu bản chất Transaction:} Đơn vị công việc logic (Unit of Work) và ý nghĩa trong toàn vẹn dữ liệu.
    \item \textbf{Làm chủ 4 tính chất ACID:} Atomicity, Consistency, Isolation, Durability và cách RDBMS hiện thực.
    \item \textbf{Sử dụng thành thạo câu lệnh:} \code{START TRANSACTION}, \code{COMMIT}, \code{ROLLBACK}, \code{SAVEPOINT}.
    \item \textbf{Kiểm soát Autocommit:} Cơ chế \code{@@autocommit} trong MySQL và quy tắc session an toàn.
    \item \textbf{Nhận biết Bẫy Implicit Commit:} Các câu lệnh DDL tự động commit ngầm và cách phòng tránh.
    \item \textbf{Nắm vững 4 mức độ Cô lập (Isolation Levels):} Ngăn chặn Dirty Read, Non-repeatable Read, Phantom Read.
    \item \textbf{Khóa \& Deadlock trong InnoDB:} Cơ chế phát hiện và giải quyết khóa chết trong môi trường đồng thời.
  \end{itemize}
\end{frame}

% Slide 3: Khai niem Transaction
\begin{frame}{1. Khái niệm Giao dịch (Transaction) trong CSDL}
  \begin{block}{Định nghĩa Giao dịch (Transaction)}
    Một \textbf{Transaction} là một chuỗi gồm một hoặc nhiều câu lệnh SQL được thực thi như \textbf{một đơn vị công việc thống nhất (Single Logical Unit of Work)}.
  \end{block}
  \vspace{0.3em}
  \begin{columns}[T]
    \begin{column}{0.48\textwidth}
      \textbf{Nguyên tắc cốt lõi:}
      \begin{itemize}
        \item \textbf{Hoặc là:} Tất cả thao tác thành công và kết quả được lưu vĩnh viễn (\code{COMMIT}).
        \item \textbf{Hoặc là:} Khi gặp sự cố, toàn bộ trạng thái quay về y hệt trước khi bắt đầu (\code{ROLLBACK}).
      \end{itemize}
    \end{column}
    \begin{column}{0.48\textwidth}
      \begin{alertblock}{Rủi ro khi thiếu Transaction}
        Nếu hệ thống sập nguồn hoặc mất mạng giữa chuỗi thao tác:
        \begin{itemize}
          \item Dữ liệu bị phân mảnh, mâu thuẫn.
          \item Vi phạm các quy tắc nghiệp vụ kinh tế.
        \end{itemize}
      \end{alertblock}
    \end{column}
  \end{columns}
\end{frame}

% Slide 4: Bai toan chuyen khoan ngan hang
\begin{frame}[fragile]{Bài toán Kinh điển: Chuyển khoản Ngân hàng}
  \textbf{Nghiệp vụ:} Chuyển 200.000 VNĐ từ Tài khoản A sang Tài khoản B:
  \begin{center}
  \begin{tikzpicture}[node distance=1.2cm, auto, >=Stealth]
    \node [rectangle, draw=dbblue, fill=dblight, rounded corners, minimum width=2.5cm, minimum height=0.8cm] (step1) {1. Kiểm tra số dư A};
    \node [rectangle, draw=dborange, fill=dborange!10, rounded corners, minimum width=2.8cm, minimum height=0.8cm, right=0.8cm of step1] (step2) {2. Trừ 200k ở TK A};
    \node [rectangle, draw=dbgreen, fill=dbgreen!10, rounded corners, minimum width=2.8cm, minimum height=0.8cm, right=0.8cm of step2] (step3) {3. Cộng 200k ở TK B};
    \node [rectangle, draw=dbblue, fill=dblight, rounded corners, minimum width=2.2cm, minimum height=0.8cm, right=0.8cm of step3] (step4) {4. Ghi Audit Log};

    \draw [->, thick] (step1) -- (step2);
    \draw [->, thick] (step2) -- (step3);
    \draw [->, thick] (step3) -- (step4);
  \end{tikzpicture}
  \end{center}
  \vspace{-0.5em}
  \begin{exampleblock}{Khi có Transaction bảo vệ}
    Nếu máy chủ gặp sự cố ở Bước 2 $\rightarrow$ Transaction tự động \textbf{ROLLBACK}, tiền của A không bị mất!
  \end{exampleblock}
\end{frame}

% Slide 5: Bo 4 tinh chat ACID
\begin{frame}{2. Bốn Tính chất Vàng: ACID}
  \begin{columns}[T]
    \begin{column}{0.48\textwidth}
      \begin{block}{\textbf{A} --- Atomicity (Tính Nguyên tử)}
        \textbf{All or Nothing:} Hoặc toàn bộ câu lệnh đều thành công, hoặc không câu lệnh nào có hiệu lực.
      \end{block}
      \vspace{0.3em}
      \begin{block}{\textbf{C} --- Consistency (Tính Nhất quán)}
        Dữ liệu luôn thỏa mãn mọi ràng buộc toàn vẹn (PK, FK, CHECK, NOT NULL) trước và sau giao dịch.
      \end{block}
    \end{column}
    \begin{column}{0.48\textwidth}
      \begin{block}{\textbf{I} --- Isolation (Tính Cô lập)}
        Các giao dịch thực thi đồng thời không can thiệp, không nhìn thấy dữ liệu chưa commit của nhau.
      \end{block}
      \vspace{0.3em}
      \begin{block}{\textbf{D} --- Durability (Tính Bền vững)}
        Sau khi đã COMMIT thành công, dữ liệu được ghi nhận vĩnh viễn vào đĩa, kể cả mất điện ngay sau đó.
      \end{block}
    \end{column}
  \end{columns}
\end{frame}

% Slide 6: Co che dam bao ACID trong MySQL InnoDB
\begin{frame}{Cách MySQL InnoDB Đảm bảo ACID}
  \begin{table}
    \centering
    \small
    \begin{tabularx}{\textwidth}{l p{4.2cm} X}
      \toprule
      \textbf{Tính chất} & \textbf{Cơ chế Kỹ thuật} & \textbf{Nguyên lý Hoạt động} \\
      \midrule
      \textbf{Atomicity} & \textbf{Undo Log} & Ghi lại trạng thái cũ của dòng dữ liệu để đảo ngược khi ROLLBACK hoặc sập máy. \\
      \textbf{Consistency} & \textbf{Schema Constraints \& Doublewrite Buffer} & Bộ máy kiểm tra khóa chính, khóa ngoại, kiểu dữ liệu, bảo vệ trang ghi dở (torn page). \\
      \textbf{Isolation} & \textbf{Locking \& MVCC} & Khóa dòng (Row locks) kết hợp Đọc đa phiên bản (Multi-Version Concurrency Control). \\
      \textbf{Durability} & \textbf{Redo Log \& WAL} & Write-Ahead Logging: Ghi nhật ký phục hồi trước khi lưu vào đĩa dữ liệu thực. \\
      \bottomrule
    \end{tabularx}
  \end{table}
\end{frame}

% Slide 7: So sanh Storage Engine InnoDB vs MyISAM
\begin{frame}{3. Storage Engine: InnoDB vs MyISAM trong MySQL}
  \begin{table}
    \centering
    \small
    \begin{tabularx}{0.92\textwidth}{l c c}
      \toprule
      \textbf{Tiêu chí Đánh giá} & \textbf{InnoDB (Khuyên dùng)} & \textbf{MyISAM (Lạc hậu)} \\
      \midrule
      Hỗ trợ Transaction & \textcolor{dbgreen}{\textbf{CÓ (Chuẩn ACID)}} & \textcolor{dbred}{\textbf{KHÔNG HỖ TRỢ}} \\
      Khả năng ROLLBACK & \textcolor{dbgreen}{\textbf{Hoàn hảo (Undo Log)}} & \textcolor{dbred}{\textbf{Không thể đảo ngược}} \\
      Cấp độ Khóa (Locking) & \textbf{Khóa dòng (Row-level)} & \textbf{Khóa toàn bảng (Table-level)} \\
      Khóa ngoại (Foreign Key) & \textcolor{dbgreen}{\textbf{Có hỗ trợ kiểm tra}} & \textcolor{dbred}{\textbf{Bỏ qua ràng buộc}} \\
      Tự phục hồi sau mất điện & \textbf{Tự động (Crash Recovery)} & \textbf{Dễ hỏng cấu trúc (Corrupt)} \\
      \bottomrule
    \end{tabularx}
  \end{table}
  \begin{alertblock}{Quy tắc vàng}
    Luôn khai báo \code{ENGINE=InnoDB} cho tất cả các bảng nghiệp vụ cần quản lý giao dịch.
  \end{alertblock}
\end{frame}

% Slide 8: Vong doi trang thai cua Transaction
\begin{frame}{4. Vòng đời Trạng thái của Transaction}
  \begin{center}
  \begin{tikzpicture}[node distance=1.6cm and 1.8cm, auto, >=Stealth]
    \node [rectangle, draw=dbblue, fill=dblight, rounded corners, minimum width=2.4cm, minimum height=0.9cm] (active) {\textbf{Active}};
    \node [rectangle, draw=dborange, fill=dborange!15, rounded corners, minimum width=3.2cm, minimum height=0.9cm, above right=0.5cm and 1.2cm of active] (pcommit) {\textbf{Partially Committed}};
    \node [rectangle, draw=dbgreen, fill=dbgreen!20, rounded corners, minimum width=2.6cm, minimum height=0.9cm, right=1.6cm of pcommit] (committed) {\textbf{Committed}};
    \node [rectangle, draw=dbred, fill=dbred!15, rounded corners, minimum width=2.4cm, minimum height=0.9cm, below right=0.5cm and 1.2cm of active] (failed) {\textbf{Failed}};
    \node [rectangle, draw=gray, fill=gray!20, rounded corners, minimum width=2.4cm, minimum height=0.9cm, right=1.6cm of failed] (aborted) {\textbf{Aborted}};

    \draw [->, thick] (active) to [bend left=15] node [above, sloped] {\scriptsize Câu lệnh DML cuối} (pcommit);
    \draw [->, thick] (pcommit) -- node [above] {\scriptsize Ghi đĩa thành công} (committed);
    \draw [->, thick] (active) to [bend right=15] node [below, sloped] {\scriptsize Lỗi hoặc Rollback} (failed);
    \draw [->, thick] (pcommit) -- node [right] {\scriptsize Ghi đĩa thất bại} (failed);
    \draw [->, thick] (failed) -- node [below] {\scriptsize Hoàn tác xong} (aborted);
  \end{tikzpicture}
  \end{center}
  \vspace{-0.5em}
  \begin{itemize}
    \item \textbf{Active:} Giao dịch đang hoạt động, thực hiện các thao tác \code{INSERT}, \code{UPDATE}, \code{DELETE}.
    \item \textbf{Partially Committed:} Thao tác cuối đã chạy trong bộ đệm, chuẩn bị ghi Redo Log xuống đĩa.
  \end{itemize}
\end{frame}

% Slide 9: Co che Autocommit
\begin{frame}[fragile]{5. Cơ chế Autocommit trong MySQL}
  \begin{block}{Biến hệ thống \code{@@autocommit}}
    Trong MySQL, mặc định \code{@@autocommit = 1}. Nghĩa là: \textbf{Mỗi câu lệnh đơn lẻ được tự động commit ngay lập tức} sau khi thực thi!
  \end{block}
  \vspace{0.4em}
  \begin{columns}[T]
    \begin{column}{0.48\textwidth}
      \textbf{Kiểm tra trạng thái autocommit:}
\begin{lstlisting}
SELECT @@autocommit;
-- Tra ve 1 (Bat) hoac 0 (Tat)
\end{lstlisting}
      \textbf{Tắt autocommit trong session:}
\begin{lstlisting}
SET autocommit = 0;
-- Moi lenh DML tiep theo se mo transaction ngam!
\end{lstlisting}
    \end{column}
    \begin{column}{0.48\textwidth}
      \textbf{Khuyến nghị thực hành:}
      \begin{itemize}
        \item Giữ \code{autocommit = 1}.
        \item Khi cần gom nhóm nghiệp vụ, mở tường minh bằng:
      \end{itemize}
\begin{lstlisting}
START TRANSACTION;
-- Thuc hien chuoi DML...
COMMIT; -- hoac ROLLBACK;
\end{lstlisting}
    \end{column}
  \end{columns}
\end{frame}

% Slide 10: Bo ba START TRANSACTION, COMMIT, ROLLBACK
\begin{frame}[fragile]{6. Bộ Ba Lệnh Cốt lõi: START, COMMIT, ROLLBACK}
  \begin{columns}[T]
    \begin{column}{0.48\textwidth}
      \begin{exampleblock}{Mẫu Transaction Thành công (Commit)}
\begin{lstlisting}
START TRANSACTION;

UPDATE bank_accounts
SET balance = balance - 200
WHERE account_id = 1;

UPDATE bank_accounts
SET balance = balance + 200
WHERE account_id = 2;

-- Xac nhan vinh vien
COMMIT;
\end{lstlisting}
      \end{exampleblock}
    \end{column}
    \begin{column}{0.48\textwidth}
      \begin{alertblock}{Mẫu Transaction Hoàn tác (Rollback)}
\begin{lstlisting}
START TRANSACTION;

UPDATE bank_accounts
SET balance = balance - 500
WHERE account_id = 1;

-- Kiem tra thay so du bi am!
-- Hoan tac sach se moi thay doi:
ROLLBACK;
\end{lstlisting}
      \end{alertblock}
    \end{column}
  \end{columns}
\end{frame}

% Slide 11: SAVEPOINT va Hoan tac tung phan
\begin{frame}[fragile]{7. Quản lý Điểm lưu với SAVEPOINT}
  \begin{block}{Mục đích của SAVEPOINT}
    Cho phép chia nhỏ transaction thành nhiều giai đoạn và \textbf{hoàn tác một phần (Partial Rollback)} mà không làm hủy toàn bộ transaction!
  \end{block}
  \vspace{0.2em}
\begin{lstlisting}
START TRANSACTION;
UPDATE bank_accounts SET balance = balance + 100 WHERE account_id = 1; -- Alice +100
SAVEPOINT sp_step1; -- Danh dau moc luu 1

UPDATE bank_accounts SET balance = balance - 9999 WHERE account_id = 2; -- Lenh loi!

-- Chi hoan tac lenh sai cua Bob, giu nguyen Alice:
ROLLBACK TO SAVEPOINT sp_step1;

-- Giai phong diem luu khoi bo nho (khong lam thay doi du lieu):
RELEASE SAVEPOINT sp_step1;

COMMIT; -- Alice van duoc cong 100 thanh cong!
\end{lstlisting}
\end{frame}

% Slide 12: Bay DDL Implicit Commit
\begin{frame}[fragile]{8. Cảnh báo Tử thần: Bẫy DDL Implicit Commit}
  \begin{alertblock}{Hiện tượng Implicit Commit trong MySQL}
    Hầu hết các câu lệnh DDL (\code{CREATE}, \code{ALTER}, \code{DROP}, \code{TRUNCATE}) \textbf{không thể rollback}. Khi gặp DDL, MySQL sẽ \textbf{tự động COMMIT ngầm} transaction đang mở trước khi chạy DDL!
  \end{alertblock}
  \vspace{0.2em}
  \begin{columns}[T]
    \begin{column}{0.5\textwidth}
      \textbf{Kịch bản sai lầm tai hại:}
\begin{lstlisting}
START TRANSACTION;
UPDATE bank_accounts SET balance = 0; -- Nguy hiem!

-- Vo tinh chay DDL:
CREATE TABLE temp_log (id INT); -- IMPLICIT COMMIT!

ROLLBACK; -- VO HIEU! Du lieu da bi luu vinh vien!
\end{lstlisting}
    \end{column}
    \begin{column}{0.46\textwidth}
      \textbf{Các lệnh gây Implicit Commit:}
      \begin{itemize}
        \item \code{CREATE TABLE}, \code{ALTER TABLE}, \code{DROP TABLE}
        \item \code{TRUNCATE TABLE}
        \item \code{LOCK TABLES}, \code{UNLOCK TABLES}
        \item \code{START TRANSACTION} (khi đang mở tx)
        \item \code{CREATE DATABASE}, \code{DROP DATABASE}
      \end{itemize}
    \end{column}
  \end{columns}
\end{frame}

% Slide 13: Cac hien tuong Concurrency
\begin{frame}{9. Các Hiện tượng Tranh chấp Đồng thời (Concurrency Anomalies)}
  Khi nhiều giao dịch chạy song song mà không cô lập đúng mức:
  \begin{columns}[T]
    \begin{column}{0.32\textwidth}
      \begin{block}{Dirty Read (Đọc rác)}
        Tx A đọc dữ liệu vừa sửa bởi Tx B \textbf{chưa commit}. Sau đó Tx B ROLLBACK $\rightarrow$ Dữ liệu A vừa đọc là hoàn toàn hư ảo!
      \end{block}
    \end{column}
    \begin{column}{0.32\textwidth}
      \begin{block}{Non-Repeatable Read}
        Tx A đọc 1 dòng, Tx B sửa dòng đó và COMMIT. Tx A đọc lại lần 2 thấy \textbf{giá trị đã bị thay đổi}!
      \end{block}
    \end{column}
    \begin{column}{0.32\textwidth}
      \begin{block}{Phantom Read (Bóng ma)}
        Tx A chạy SELECT WHERE, Tx B chèn thêm dòng mới thỏa điều kiện. Tx A SELECT lại thấy \textbf{xuất hiện thêm dòng bóng ma}!
      \end{block}
    \end{column}
  \end{columns}
\end{frame}

% Slide 14: Bon muc do co lap (Isolation Levels)
\begin{frame}{10. Bốn Mức độ Cô lập (Isolation Levels)}
  \begin{table}
    \centering
    \scriptsize
    \begin{tabularx}{\textwidth}{l c c c X}
      \toprule
      \textbf{Mức độ Cô lập} & \textbf{Dirty Read} & \textbf{Non-Repeatable} & \textbf{Phantom} & \textbf{Đặc điểm \& Ứng dụng} \\
      \midrule
      \textbf{READ UNCOMMITTED} & \textcolor{dbred}{Bị lỗi} & \textcolor{dbred}{Bị lỗi} & \textcolor{dbred}{Bị lỗi} & Nhanh nhất, không an toàn. \\
      \textbf{READ COMMITTED} & \textcolor{dbgreen}{Ngăn chặn} & \textcolor{dbred}{Bị lỗi} & \textcolor{dbred}{Bị lỗi} & Mặc định Oracle, PostgreSQL. \\
      \textbf{REPEATABLE READ} & \textcolor{dbgreen}{Ngăn chặn} & \textcolor{dbgreen}{Ngăn chặn} & \textcolor{dbgreen}{Ngăn chặn*} & \textcolor{dbblue}{\textbf{MẶC ĐỊNH MYSQL INNODB}}. Dùng Next-Key Locks. \\
      \textbf{SERIALIZABLE} & \textcolor{dbgreen}{Ngăn chặn} & \textcolor{dbgreen}{Ngăn chặn} & \textcolor{dbgreen}{Ngăn chặn} & Tuần tự hóa, hiệu năng thấp nhất. \\
      \bottomrule
    \end{tabularx}
  \end{table}
  \vspace{0.4em}
  \textbf{Cú pháp thiết lập trong MySQL:}
  \begin{center}
    \code{SET SESSION TRANSACTION ISOLATION LEVEL READ COMMITTED;}
  \end{center}
\end{frame}

% Slide 15: Co che Khoa (Locking)
\begin{frame}{11. Cơ chế Khóa (Locking) trong InnoDB}
  \begin{columns}[T]
    \begin{column}{0.48\textwidth}
      \begin{block}{Khóa Chia sẻ --- Shared Lock (S)}
        \begin{itemize}
          \item Cho phép nhiều transaction cùng \textbf{ĐỌC (Read)}.
          \item Ngăn không cho transaction khác ghi/sửa.
          \item Cú pháp: \code{SELECT ... FOR SHARE;} (hoặc \code{LOCK IN SHARE MODE}).
        \end{itemize}
      \end{block}
    \end{column}
    \begin{column}{0.48\textwidth}
      \begin{alertblock}{Khóa Độc quyền --- Exclusive Lock (X)}
        \begin{itemize}
          \item Cấp cho transaction muốn \textbf{GHI (Write)}.
          \item Ngăn cả đọc lẫn ghi của tất cả transaction khác.
          \item Tự động cấp khi chạy \code{UPDATE}, \code{DELETE}, hoặc chủ động: \code{SELECT ... FOR UPDATE;}.
        \end{itemize}
      \end{alertblock}
    \end{column}
  \end{columns}
\end{frame}

% Slide 16: Hien tuong Deadlock
\begin{frame}[fragile]{12. Hiện tượng Deadlock (Khóa Chết) \& Xử lý}
  \begin{columns}[T]
    \begin{column}{0.5\textwidth}
      \textbf{Kịch bản Deadlock điển hình:}
      \begin{itemize}
        \item \textbf{Tx 1:} Giữ khóa dòng A, đang chờ khóa dòng B.
        \item \textbf{Tx 2:} Giữ khóa dòng B, đang chờ khóa dòng A.
        \item $\rightarrow$ \textbf{Vòng tròn phụ thuộc vĩnh viễn!}
      \end{itemize}
      \vspace{0.3em}
      \begin{exampleblock}{InnoDB Deadlock Detection}
        InnoDB tự động phát hiện đồ thị chờ đợi, chọn transaction có khối lượng thao tác nhỏ hơn làm \textbf{nạn nhân để ROLLBACK}, báo lỗi:
        \code{ERROR 1213: Deadlock found}.
      \end{exampleblock}
    \end{column}
    \begin{column}{0.46\textwidth}
      \textbf{Quy tắc phòng ngừa Deadlock:}
      \begin{enumerate}
        \item \textbf{Cố định thứ tự truy cập:} Luôn khóa các bảng/dòng theo cùng một trình tự logic (VD: sắp xếp tăng dần theo ID).
        \item \textbf{Giữ transaction ngắn:} Chỉ thực hiện các thao tác DML cần thiết, commit càng sớm càng tốt.
        \item \textbf{Tạo chỉ mục (Index) đầy đủ:} Giúp InnoDB khóa chính xác từng dòng thay vì khóa quét cả bảng!
      \end{enumerate}
    \end{column}
  \end{columns}
\end{frame}

% Slide 17: Quy trinh an toan khi viet Transaction
\begin{frame}{13. Quy trình An toàn khi Xử lý Transaction trong Ứng dụng}
  \begin{center}
  \begin{tikzpicture}[node distance=1.1cm, auto, >=Stealth]
    \node [rectangle, draw=dbblue, fill=dblight, rounded corners, minimum width=9cm, minimum height=0.7cm] (b1) {1. Bắt đầu Transaction (\code{START TRANSACTION})};
    \node [rectangle, draw=dbblue, fill=dblight, rounded corners, minimum width=9cm, minimum height=0.7cm, below=0.3cm of b1] (b2) {2. Thực thi chuỗi câu lệnh DML trong khối Try--Catch};
    \node [rectangle, draw=dborange, fill=dborange!15, rounded corners, minimum width=9cm, minimum height=0.7cm, below=0.3cm of b2] (b3) {3. Kiểm tra tính hợp lệ \& Nghiệp vụ (Số dư, Tồn kho, Trạng thái)};
    \node [rectangle, draw=dbgreen, fill=dbgreen!20, rounded corners, minimum width=4.3cm, minimum height=0.8cm, below left=0.4cm and -0.2cm of b3] (b4a) {HỢP LỆ: \textbf{COMMIT}};
    \node [rectangle, draw=dbred, fill=dbred!20, rounded corners, minimum width=4.3cm, minimum height=0.8cm, below right=0.4cm and -0.2cm of b3] (b4b) {CÓ LỖI: \textbf{ROLLBACK}};

    \draw [->, thick] (b1) -- (b2);
    \draw [->, thick] (b2) -- (b3);
    \draw [->, thick] (b3) -| (b4a);
    \draw [->, thick] (b3) -| (b4b);
  \end{tikzpicture}
  \end{center}
  \vspace{-0.5em}
  \begin{alertblock}{Nguyên tắc tối hậu}
    Tuyệt đối không để một transaction mở lơ lửng mà không có \code{COMMIT} hoặc \code{ROLLBACK} dứt điểm!
  \end{alertblock}
\end{frame}

% Slide 18: Huong dan thuc hanh Lab tuong tac
\begin{frame}{14. Thực hành Lab Tương tác trên Trình duyệt}
  \begin{columns}[T]
    \begin{column}{0.5\textwidth}
      \textbf{Tài liệu \& Lab Thực hành Tương tác:}
      \begin{itemize}
        \item File: \code{transaction.html}
        \item Chạy trực tiếp 100\% trên trình duyệt (không cần cài đặt server).
        \item Bộ mô phỏng Dual-Engine: SQLite WebAssembly + JS Fallback.
      \end{itemize}
      \vspace{0.3em}
      \textbf{Tính năng nổi bật:}
      \begin{itemize}
        \item Theo dõi biến động số dư và cờ Uncommitted theo thời gian thực.
        \item Trực quan hóa Savepoints stack.
      \end{itemize}
    \end{column}
    \begin{column}{0.46\textwidth}
      \begin{exampleblock}{Hệ thống Tự chấm điểm}
        \begin{itemize}
          \item \textbf{15 Bài tập SQL tự chấm} từ cơ bản đến phức tạp (Chuyển khoản, E-Commerce, Savepoint).
          \item \textbf{10 Câu hỏi trắc nghiệm ACID} củng cố kiến thức lý thuyết.
          \item \textbf{Xuất chứng chỉ PNG 2x High-DPI} kèm điểm số và mã xác thực học viên.
        \end{itemize}
      \end{exampleblock}
    \end{column}
  \end{columns}
\end{frame}

% Slide 19: Tong ket bai hoc
\begin{frame}{15. Tổng kết Bài học (Key Takeaways)}
  \begin{enumerate}
    \item \textbf{Transaction} là đơn vị công việc logic cốt lõi đảm bảo dữ liệu không bao giờ bị dở dang.
    \item Bốn tính chất \textbf{ACID} là nền tảng của hệ quản trị cơ sở dữ liệu quan hệ (RDBMS).
    \item Trong MySQL, luôn ưu tiên sử dụng storage engine \textbf{InnoDB} để được hỗ trợ transaction đầy đủ.
    \item Hiểu rõ cơ chế \textbf{Autocommit}: Mặc định bật, cần dùng \code{START TRANSACTION} để gom nhóm DML.
    \item Sử dụng \textbf{SAVEPOINT} khi cần chia nhỏ giai đoạn và hoàn tác một phần.
    \item Hết sức cảnh giác với \textbf{Implicit Commit}: Các câu lệnh DDL sẽ ngầm lưu vĩnh viễn các thao tác trước đó.
    \item Nắm vững 4 mức \textbf{Isolation}: MySQL InnoDB mặc định dùng \textbf{REPEATABLE READ}.
  \end{enumerate}
\end{frame}

% Slide 20: Hoi dap & Ket thuc
\begin{frame}
  \begin{center}
    \vspace{1.5em}
    {\huge \textcolor{dbblue}{\textbf{Hỏi \& Đáp (Q\&A)}}}\\[1.2em]
    {\large Cảm ơn các bạn đã chú ý lắng nghe!}\\[1.5em]
    \textbf{Tài nguyên học tập:}\\[0.5em]
    Trang Web Khóa học: \href{https://vudmvn.github.io/dbms-course}{\code{vudmvn.github.io/dbms-course}}\\[0.3em]
    Lab Tương tác: \code{lectures/MySQL/data-modification-transactions/transaction.html}
  \end{center}
\end{frame}

\end{document}
